\chapter{Using the Code}

The repository is available at \href{\projecturl}{\projecturl}. Python 3.10 or
newer, NumPy, ngspice, and a compatible SKY130A ngspice model library are
required for real-circuit evaluation. Commands below assume the repository root
as the working directory.

\section{Create an environment}

\begin{lstlisting}[language=bash]
python -m venv .venv
# Windows PowerShell
.\.venv\Scripts\Activate.ps1
python -m pip install -r requirements.txt
\end{lstlisting}

Set the external tools without committing machine-specific paths:

\begin{lstlisting}[language=bash]
$env:NGSPICE_EXECUTABLE = "C:\path\to\ngspice_con.exe"
$env:SKY130_MODEL_LIBRARY = "C:\path\to\sky130.lib.spice"
\end{lstlisting}

\section{Run the regression suite}

\begin{lstlisting}[language=bash]
python -m unittest discover -v
\end{lstlisting}

Tests requiring the SKY130 model are skipped when the model variable is not
configured. The exact executed, passed, failed, and skipped counts should be
recorded together with the date, commit, and interpreter.

\section{Evaluate one receiver design}

\begin{lstlisting}[language=bash]
python -m experiments.evaluate_receiver --fidelity training
\end{lstlisting}

To use the included public development channel:

\begin{lstlisting}[language=bash]
python -m experiments.evaluate_receiver \
  --channel channels/ieee802_ibm_20db_thru.s4p \
  --channel-ports 1 3 2 4 --fidelity candidate
\end{lstlisting}

The port order is one-based and represents TXP, TXN, RXP, and RXN. It should be
overridden only from documented channel-port information.

\section{Train or run the optimization pipeline}

Inspect the complete command-line options before starting an expensive run:

\begin{lstlisting}[language=bash]
python -m experiments.run_autockt_pipeline --help
python -m experiments.train_rl --help
\end{lstlisting}

Use a synthetic backend for a fast software-path check and the real backend for
electrical evidence. Training and inference must be distinguished: training
updates policy weights, while inference loads a fixed checkpoint and generates
candidates.

\section{Launch the local interface}

\begin{lstlisting}[language=bash]
python experiments/web_ui.py --port 8001
\end{lstlisting}

Open \url{http://127.0.0.1:8001}. Choose the backend, target, checkpoint,
channel, PVT scope, and optional final measurements. ``None'' is nominal-only;
``smoke'' is a diagnostic and not a robustness proof; larger PVT sets can take
hours. Preserve the generated JSON, event log, graph, schematic, and eye
artifact for every result quoted in a report.

The reviewed UI accepts four structured numeric target fields. It is not a
natural-language interface. If an LLM wrapper is added later, it should only map
text to a validated \code{TargetSpec} and pipeline configuration; all electrical
verdicts must continue to come from the simulator and deterministic assessment
code.

\section{Build this report}

From the \code{report/} directory:
\begin{lstlisting}[language=bash]
pdflatex main.tex
bibtex main
pdflatex main.tex
pdflatex main.tex
\end{lstlisting}
The same source can be uploaded to Overleaf. The two supplied CTLE/receiver PDFs
are retained as vector figures, so no raster conversion is required.
